\documentclass{notaaula}

        \titulonota{Potência e consumo elétrico}
        \creditos{Turma 3A · Física · Setembro/2026 · Prof. Josué Cavalcante}

        \begin{document}
        \cabecalho

        \begin{sobre}
        Eletrodinâmica: potência elétrica, consumo de energia, quilowatt-hora e rendimento. Habilidades em foco: EM13CNT107, EM13CNT203 e EM13CNT304.
        \end{sobre}

        \section{Potência elétrica}

\begin{copiar}{Taxa de transformação}
A potência indica a rapidez com que a energia elétrica é transformada:
\[
  P=Ui.
\]
Para resistor ôhmico, usando $U=Ri$,
\[
  P=Ri^2=\frac{U^2}{R}.
\]
\simbolos{$P$ (W); $U$ (V); $i$ (A); $R$ ($\Omega$).}
\atencao{Escolha a expressão compatível com as grandezas conhecidas; não misture valores de aparelhos diferentes.}
\end{copiar}

\begin{exemplo}
Um chuveiro opera em $220\un{V}$ e corrente $25\un{A}$.
\[
  P=Ui=220\cdot25=5500\un{W}=
  \resultado{\dec{5,5}\un{kW}}.
\]
Sua resistência de operação é $R=U/i=\dec{8,8}\,\Omega$.
\end{exemplo}

\section{Energia e quilowatt-hora}

\begin{copiar}{Consumo}
Para potência constante,
\[
  E=P\Delta t.
\]
A unidade do SI é o joule, mas a conta de eletricidade usa quilowatt-hora:
\[
  1\un{kWh}=3{,}6\times10^6\un{J}.
\]
Para obter kWh, use potência em kW e tempo em horas. O custo simples é
\[
  C=E_{\mathrm{kWh}}\times\text{tarifa}.
\]
\end{copiar}

\begin{exemplo}[Consumo mensal]
O chuveiro de $\dec{5,5}\un{kW}$ é usado $20$ minutos por dia durante 30 dias.
O tempo total é $10\un{h}$; portanto,
\[
  E=\dec{5,5}\cdot10=\resultado{55\un{kWh}}.
\]
Com tarifa de $R\$\,\dec{0,90}/\un{kWh}$, o custo é $R\$\,\dec{49,50}$.
\end{exemplo}

\section{Rendimento e uso consciente}

\begin{copiar}{Energia útil e perdas}
\[
  \eta=\frac{E_{\mathrm{útil}}}{E_{\mathrm{recebida}}}
  =\frac{P_{\mathrm{útil}}}{P_{\mathrm{recebida}}}.
\]
Nenhum aparelho real tem rendimento acima de $100\%$. Reduzir desperdícios envolve diminuir o
tempo de uso, escolher tecnologia mais eficiente e adequar a potência à necessidade.
\diaadia{Em lâmpadas, a parcela luminosa é útil; aquecimento indesejado representa perda para essa finalidade.}
\end{copiar}

\begin{exemplo}[Rendimento]
Um motor recebe $800\un{W}$ e fornece $600\un{W}$ de potência mecânica.
\[
  \eta=\frac{600}{800}=\dec{0,75}=
  \resultado{75\%}.
\]
As perdas somam $200\un{W}$.
\end{exemplo}

\section{Segurança e leitura da conta}

\begin{copiar}{Corrente e proteção}
A corrente total de aparelhos ligados em paralelo soma-se. Condutores e disjuntores devem ser
dimensionados para essa corrente. Sobrecarga eleva o aquecimento por efeito Joule e aumenta o risco
de incêndio. A conta registra energia, não potência: kWh não é kW/h.
\end{copiar}

\begin{dica}
Um aparelho potente pode consumir pouca energia se funcionar por pouco tempo. Consumo depende de
potência e duração: compare sempre $P\Delta t$.
\end{dica}

\begin{exercicios}
\nivel{1}{Conceitos}
\begin{questoes}
\item Calcule a potência de aparelho em $120\un{V}$ percorrido por $5\un{A}$.
\item Converta $2\un{kWh}$ em joules.
\item Diferencie potência de energia elétrica.
\item Um resistor de $20\,\Omega$ é percorrido por $2\un{A}$. Ache $P$.
\item O que significa rendimento de $80\%$?
\nivel{2}{Aplicação}
\item Uma TV de $150\un{W}$ funciona $4\un{h}$. Calcule o consumo em kWh.
\item Um ferro de $1\un{kW}$ funciona 30 minutos por dia durante 20 dias. Ache o consumo.
\item Com tarifa de $R\$\,\dec{0,80}/\un{kWh}$, calcule o custo do item anterior.
\item Um motor recebe $2\un{kW}$ e tem rendimento de $70\%$. Ache a potência útil.
\nivel{3}{Síntese}
\item Dois chuveiros têm $5500\un{W}$ e $6800\un{W}$. Para o mesmo tempo de uso, qual consome mais e em que razão?
\item Explique por que vários aparelhos potentes na mesma tomada podem causar sobrecarga.
\item Proponha duas ações mensuráveis para reduzir o consumo mensal de energia.
\end{questoes}
\gabarito{1) $600\un{W}$; 2) $7{,}2\times10^6\un{J}$; 4) $80\un{W}$;
6) $\dec{0,60}\un{kWh}$; 7) $10\un{kWh}$; 8) $R\$\,\dec{8,00}$;
9) $\dec{1,4}\un{kW}$; 10) o de $6800\un{W}$, razão $68/55\approx\dec{1,24}$.}
\end{exercicios}


\end{document}
